\chapter{Conclusioni e Sviluppi futuri}
\label{cap:conclusioni}

\section{Risultati}

Il lavoro ha risposto alla domanda posta nell'introduzione: con sole fonti pubbliche è possibile costruire uno strumento che affianca, per distretto socio-sanitario del Salento, qualità dell'aria, calore e mortalità. Il prodotto è completo dall'acquisizione dei dati alla risposta in linguaggio naturale, e gira interamente in locale. Se ne riassumono i risultati in quattro punti.
\begin{itemize}
  \item \textbf{Dati}: quattro fonti d'aria, due di temperatura, due raster di mitigazione e tre pubblicazioni sulla mortalità sono riportate a una grana comune, con stati espliciti per ciò che manca. La ridondanza fra fonti indipendenti le valida (Tabella~\ref{tab:validazione}): la mortalità ISTAT coincide con Eurostat al 100\% e ERA5-Land riproduce la temperatura ARPA con correlazione 0,97--0,99.
  \item \textbf{Copertura dichiarata}: 124 comuni su 145 non hanno stazioni di misura. Dove manca il dato misurato il sistema mostra quello modellato, etichettato come tale, invece di un buco o di un valore non supportato.
  \item \textbf{Dashboard}: un'interfaccia per decidere, con la carta, la sintesi per finestra, il confronto fra fonti e fra domini e un sommario stagionale.
  \item \textbf{Assistente}: un modello da 4 miliardi di parametri in locale, vincolato da codice a un elenco chiuso di analisi e argomenti. Nella batteria di prova nessuna risposta è uscita dal perimetro.
\end{itemize}

\section{Limiti}

I limiti sono quelli annunciati e quelli emersi lavorando.
\begin{itemize}
  \item \textbf{Esplorativo, non causale}: fumo, età, comorbidità e aria indoor non sono osservabili. L'esempio dei dati simulati (Sezione~\ref{sec:domini-correlazione}) mostra che una correlazione grezza può avere persino il segno sbagliato. Metriche formali di validazione della relazione fra i domini non sono state definite; lo strumento mostra, non conferma ipotesi.
  \item \textbf{Salute simulata}: ricoveri, accessi in pronto soccorso e assistenza domiciliare sono simulati, perché non esiste una fonte pubblica sotto la provincia. I dati reali, richiesti ad AReSS Puglia, non sono arrivati. Gli effetti di PM\textsubscript{10} e NO\textsubscript{2} imposti nella simulazione non si distinguono dal caso: è un limite di potenza che varrà anche per i dati reali. Resta la mortalità per tutte le cause, senza distinzione di causa.
  \item \textbf{Fonti}: la fonte primaria d'aria, ARPA Puglia, è raggiunta tramite un indirizzo non documentato, perché l'API ufficiale non risponde; va confermato l'uso con l'ente. Le serie di OpenAQ presentano orari sfasati, che non si correggono ma si convertono come dichiarati. Ogni dato di CAMS ha una risoluzione di circa 10\,km: per i comuni piccoli è una stima, non una misura.
  \item \textbf{Calore}: l'isola di calore in senso stretto non è calcolata; il progetto mostra l'esposizione e la mitigazione. Per il Salento la letteratura indica che l'esito dipende dai materiali e dall'uso del suolo (Sezione~\ref{sec:domini-uhi}).
  \item \textbf{Assistente}: la verifica controlla i numeri e non le parole, perciò una risposta può usare un termine sbagliato con numeri giusti. Per questo ogni risposta porta accanto i dati usati. Il perimetro è provato su 64 domande: è una prova di assenza di errori noti, non una garanzia.
\end{itemize}

\section{Sviluppi futuri}

\begin{itemize}
  \item \textbf{Isola di calore da satellite}: stimare la differenza di temperatura superficiale urbano--rurale per comune e anno da immagini Landsat Collection 2 Level-2, come quarta grandezza accanto alle altre. È lo sviluppo più coerente con il quadro del Capitolo~\ref{cap:1}, e richiede di trattare le nuvole e di definire in modo esplicito che cosa è urbano e rurale. Sentinel-3 ed ECOSTRESS sono alternative da valutare per copertura temporale e risoluzione.
  \item \textbf{Dati sanitari reali}: sostituire i dati simulati con quelli di un ente che li possa cedere. Il formato è già uguale; un'unità sanitaria potrebbe fornire anche le cause. A quel punto vale la pena definire le metriche di validazione.
  \item \textbf{Perimetro e fonti}: documentare la fonte ARPA e ottenere la conferma d'uso, e estendere i confronti ad altre aree.
  \item \textbf{Assistente}: collegare le analisi su serie mensili ancora scoperte, valutare modelli più grandi e provare la ricerca per \emph{embedding}, se il corpus crescerà.
  \item \textbf{Pubblicazione}: separare presentazione, codice e interfaccia in repository distinti, con licenza dell'azienda partner e senza dati.
\end{itemize}
